% Problemas del capítulo 13. Generado por tools/problemas/gen.py; no editar a mano.
\section{Problemas}\label{sec:problemas-13}

Graduados: \(\star\) directo, \(\star\star\) para pensar, \(\star\star\star\) desafío. Las respuestas están en el apéndice~\ref{ap:respuestas}; conviene intentarlos antes de mirarlas.

\begin{problema}[\({\star}\)]\label{pb:13-1}
¿Qué curvatura tiene una ruta recta? ¿Y una rotonda de \(20\)~m de radio?
\end{problema}

\begin{problema}[\({\star}{\star}\)]\label{pb:13-2}
Un camino sigue el perfil \(y=x^2/200\) (en metros). Calculá la curvatura y el radio en \(x=0\) y en \(x=50\).
\end{problema}

\begin{problema}[\({\star}{\star}\)]\label{pb:13-3}
Cerca de su punto más bajo, una superficie es \(z=\frac{x^2}{2\cdot50}+\frac{y^2}{2\cdot200}\). Calculá las curvaturas principales, la curvatura de Gauss y, con la fórmula de Euler, la curvatura en la dirección que forma \(30^\circ\) con el eje \(x\).
\end{problema}

\begin{problema}[\({\star}{\star}\)]\label{pb:13-4}
La silla \(z=xy/100\): calculá las curvaturas principales en el origen y sus direcciones. ¿Qué signo tiene la curvatura de Gauss?
\end{problema}

\begin{problema}[\({\star}{\star}{\star}\)]\label{pb:13-5}
Sobre una esfera de radio \(R=6371\)~km se corta con un plano que contiene una recta tangente y forma \(60^\circ\) con la normal. ¿Qué radio tiene la circunferencia del corte? (Fórmula de Meusnier.)
\end{problema}
